\chapter{Leidingdimensionering}
\label{h:leiding}

\section{Methode}
\begin{methode}[– Cengel H5, behoud van massa]
$\md = \rho\,\Vd = \rho\,A\,v$. Bij een gekozen maximale snelheid $v_{\max}$ volgt de minimale doorsnede, en daarmee
\[
\Vd = \frac{\md}{3600\,\rho}\ [\mathrm{m^3/s}],\qquad
d_{\min} = \sqrt{\frac{4\,\Vd}{\pi\,v_{\max}}},\qquad
v_\mathrm{werkelijk} = \frac{\Vd}{\frac{\pi}{4}d_i^2}.
\]
Daarna kies je de kleinste standaardpijp met $d_i \ge d_{\min}$. Hier is dat RVS volgens ASME B36.19, schedule 40S
($d_i = d_\mathrm{buiten} - 2\cdot$wanddikte).
\end{methode}
De dichtheid volgt uit de toestand van de stroom:
%% if A
CoolProp met ($T$, $P$) voor \HH\ en water, en ($T$, $x$) of ($P$, $h$) voor de koudemiddelen. Voor \OO\ met damp geldt de
ideale-gaswet voor het mengsel, $\rho = P\,\bar M/(R_u T)$.
%% else
ideaal gas $\rho = P/(RT)$ voor \HH\ en $\rho = P\,\bar M/(R_u T)$ voor het \OO-dampmengsel; $\rho = 1/v_f(T)$ uit tabel A-4 voor
water (vloeistof is bijna onsamendrukbaar, dus $v \approx v_f(T)$, Cengel H3); en voor de koudemiddelen afgelezen op de
isochoren van het log $p$–$h$-diagram.
%% endif

\section{Maximale snelheden}
\begin{center}
\begin{tabularx}{\textwidth}{lrX}
\toprule
categorie & $v_{\max}$ [m/s] & waarom\\
\midrule
\HH-gas & \VAR{n(I.v_max_H2,0)} & richtwaarde voor waterstof (EIGA Doc 121); hogere snelheden geven erosie en geluid\\
\OO-gas (droog/verzadigd) & \VAR{n(I.v_max_O2,0)} & zuurstof: deeltjes die met hoge snelheid inslaan kunnen ontsteking veroorzaken (EIGA Doc 13), dus conservatief\\
\OO\ met druppels (S14) & \VAR{n(I.v_max_nat,0)} & tweefasestroming, voorkomt meesleuren en erosie\\
vloeistof persleiding & \VAR{n(I.v_max_pers,0)} & gangbaar compromis tussen drukval en leidingkosten\\
vloeistof zuigleiding & \VAR{n(I.v_max_zuig,0)} & lage drukval aan de zuigkant (NPSH)\\
koudemiddeldamp & \VAR{n(I.v_max_ref_damp,0)} & drukval in zuig- en persleiding kost COP\\
koudemiddelvloeistof & \VAR{n(I.v_max_ref_vl,0)} & voorkomt flash door drukval vóór het smoorventiel\\
koudemiddel tweefase & \VAR{n(I.v_max_ref_2f,0)} & tussenwaarde voor het mengsel na het smoorventiel\\
\bottomrule
\end{tabularx}
\end{center}

\section{Twee voorbeelden met de hand}
\BLOCK{set L4 = R.leidingen|selectattr("id","equalto","S4")|first}
\BLOCK{set LN = R.leidingen|selectattr("id","equalto","N1")|first}
\kop{S4: \HH\ van ITM naar V-101 (per unit)}
%% if A
$\rho = $ \texttt{PropsSI("D","T",298.15,"P",\VAR{n((I.P_H2_ITM*100+PATM)*1000,0)},"Hydrogen")} $= \VAR{n(L4.rho,3)}$\,kg/m$^3$.
%% else
$\rho = \frac{P}{RT} = \frac{\VAR{n(I.P_H2_ITM*100+PATM,1)}}{4{,}124\cdot298{,}15} = \VAR{n(L4.rho,3)}$\,kg/m$^3$.
%% endif
\[
\Vd = \frac{\VAR{n(L4.m,0)}}{3600\cdot\VAR{n(L4.rho,3)}} = \VAR{n(L4.V/3600,5)}\,\mathrm{m^3/s},\quad
d_{\min} = \sqrt{\frac{4\cdot\VAR{n(L4.V/3600,5)}}{\pi\cdot\VAR{n(L4.v_max,0)}}} = \VAR{n(L4.d_min,1)}\,\mathrm{mm}
\]
\[
\Rightarrow\; \text{DN\,\VAR{L4.DN}} \ (d_i = \VAR{n(L4.d_i,2)}\,\mathrm{mm}),\ v = \VAR{n(L4.v,1)}\,\mathrm{m/s}.
\]
\kop{N1: warmtenet retour (installatie)}
$\rho = \VAR{n(LN.rho,1)}$\,kg/m$^3$ (water bij 40\degC),
\[
\Vd = \frac{\VAR{n(LN.m,0)}}{3600\cdot\VAR{n(LN.rho,1)}} = \VAR{n(LN.V/3600,4)}\,\mathrm{m^3/s},\quad
d_{\min} = \sqrt{\frac{4\cdot\VAR{n(LN.V/3600,4)}}{\pi\cdot\VAR{n(LN.v_max,0)}}} = \VAR{n(LN.d_min,0)}\,\mathrm{mm}
\]
\[
\Rightarrow\; \text{DN\,\VAR{LN.DN}},\ v = \VAR{n(LN.v,2)}\,\mathrm{m/s}.
\]

\section{Alle leidingen}
\begin{center}\scriptsize\setlength{\tabcolsep}{3pt}
\begin{longtable}{lp{5.4cm}lrrrrrrr}
\toprule
stroom & omschrijving & schaal & $\md$ [kg/h] & $\rho$ [kg/m$^3$] & $\Vd$ [m$^3$/h] & $v_{\max}$ & $d_{\min}$ [mm] & DN & $v$ [m/s]\\
\midrule\endhead
%% for L in R.leidingen
\VAR{L.id} & \VAR{tex(L.naam)} & \VAR{"unit" if L.schaal == "per unit" else "inst."} & \VAR{n(L.m, 1 if L.m < 1000 else 0)} & \VAR{g(L.rho,4)} & \VAR{g(L.V,4)} & \VAR{n(L.v_max,0)} & \VAR{n(L.d_min,1)} & \VAR{L.DN} & \VAR{n(L.v,2)}\\
%% endfor
\bottomrule
\end{longtable}
\end{center}
\begin{waarom}[ontwerppunt en kanttekeningen]
Gedimensioneerd op vollast met net en warmtepompen aan; S9 op het maximale recycledebiet (\VAR{n(I.m_H2*I.m_min_pct/100,0)}\kgh).
De zuigleiding van K-201 (R1) is groot (DN\,\VAR{(R.leidingen|selectattr("id","equalto","R1")|first).DN}), omdat de damp
bij \VAR{n(wp1.Pe,0)}\kPa\ een lage dichtheid heeft en er ruim \VAR{n(wp1.m_ref/1000,0)}\,ton/h, oftewel
\VAR{n(wp1.m_ref,0)}\kgh\ koudemiddel rondgaat. In de praktijk worden de warmtepompen in meerdere parallelle machines
uitgevoerd, met elk een eigen kleinere leiding. S10 (blowoff) dimensioneer je op het afblaasgeval, niet op het
procesdebiet.
\end{waarom}
